\subsection{Aplicación de TCP y gRPC al PFC desde CAP}
\label{sec:pfc-sgra}

El proyecto considerado es el Sistema de Gestión de Refuerzos Académicos
(SGRA), organizado en un frontend Angular y un backend Spring Boot con Java 21
y acceso a PostgreSQL. Los servicios del frontend realizan peticiones HTTP para
crear solicitudes, consultar sesiones y registrar asistencia o resultados. La
API del backend contempla aceptar, reprogramar y cancelar refuerzos, además de
gestionar sus recursos \cite{belinda_sgra_front,belinda_sgra_back}.
Esta separación permite ubicar las comunicaciones del sistema; mantener dos
repositorios no demuestra por sí solo que exista una arquitectura de
microservicios o una base de datos replicada.

\paragraph{Aplicación de sockets TCP.}
Una extensión posible sería conectar un agente de los espacios de refuerzo con
el backend para comunicar cambios de disponibilidad. El agente enviaría mensajes
con identificador de operación, espacio, franja horaria y estado, y recibiría un
acuse que permitiera relacionar la respuesta con el envío. La delimitación y
lectura completa utilizadas en la práctica servirían como base para ese protocolo
de aplicación. El acuse tendría que distinguir la recepción del mensaje de la
confirmación de una reserva, que exige validar y guardar la operación.
Este canal sería adicional a la API HTTP que utiliza el navegador.

\paragraph{Aplicación de gRPC.}
Si la programación de sesiones y la gestión de espacios se separaran en
servicios, podrían comunicarse mediante un contrato Protocol Buffers con
operaciones como consultar disponibilidad y reservar un espacio. Los mensajes
incluirían la sesión, el espacio, el horario y un identificador que permitiera
detectar solicitudes repetidas. gRPC facilitaría definir esos métodos y generar
clientes y servidores para su comunicación interna
\cite{belinda_grpc_introduccion}. Se mantendría la API HTTP para el frontend.
Esta alternativa plantea una evolución del sistema; no describe un servicio
gRPC ya incorporado a SGRA.

\paragraph{Decisiones ante una partición de red.}
CAP resulta pertinente si el estado de las reservas se distribuye y replica.
Ante una partición, dos réplicas incomunicadas podrían intentar confirmar el
mismo espacio y horario. Para preservar una única decisión válida se priorizaría
la consistencia de la reserva: solo se confirmaría donde pudiera coordinarse
la escritura con la autoridad o el cuórum definido. El componente aislado
devolvería indisponibilidad o dejaría la solicitud pendiente, sin presentarla
como una reserva confirmada. Se sacrificaría así disponibilidad de esa operación
durante la partición \cite{belinda_lee2023_cap}.

Para consultas informativas, como el historial, podría permitirse una copia
temporalmente desactualizada, indicando su fecha de actualización. Esa tolerancia
no se extendería a confirmar reservas con datos antiguos. Son decisiones por
operación y bajo una arquitectura replicada; la revisión de los repositorios
no basta para clasificar el despliegue actual de SGRA como CP o AP.

La práctica aporta mecanismos para intercambiar mensajes y conservar trazas que
ayuden a diagnosticar solicitudes y acuses. El reloj Lamport permite relacionar
su orden causal, pero no evita por sí solo una doble reserva. Tampoco elegir TCP
o gRPC resuelve CAP: las garantías dependen de cómo el servicio coordine y
almacene su estado. Los RTT obtenidos en localhost deberán interpretarse dentro
de sus condiciones de prueba y no como una predicción directa del rendimiento
de SGRA en producción.
